\documentclass[preprint,11pt]{elsarticle}

\usepackage{amsmath}
\usepackage{amssymb}
\usepackage{mathtools}
\usepackage{amsthm}
\usepackage{booktabs}
\usepackage[hidelinks]{hyperref}

\newtheorem{lemma}{Lemma}
\newtheorem{proposition}{Proposition}
\newtheorem{remark}{Remark}

% --- convenience macros ---
\newcommand{\cs}{c_s}
\newcommand{\eq}{\mathrm{eq}}
\newcommand{\dd}{\,\mathrm{d}}
\newcommand{\Deltafour}{\Delta_{\alpha\beta\gamma\delta}}
\newcommand{\sumi}{\sum_i}

\journal{Notes on the lattice Boltzmann equilibrium moment hierarchy}

\begin{document}

\begin{frontmatter}

\title{On the velocity moments of the lattice Boltzmann equilibrium:
second- and third-order results, and the structural cubic defect of the
D3Q27 lattice}

\author{Peng-Chung Chen}
\address{Computational Fluid Dynamics Laboratory, Department of Power
Mechanical Engineering, National Tsing Hua University, Hsinchu 30013, Taiwan}

\begin{abstract}
The discrete velocity moments of the lattice Boltzmann (LB) equilibrium are
frequently quoted in two superficially contradictory forms in the
literature: a kinetic-theory form that carries the cubic velocity term
$\rho\,u_\alpha u_\beta u_\gamma$ at third order, and a discrete form that does
not. We resolve this apparent contradiction by deriving, from first
principles and in full, four distinct objects: the second- and third-order
moments of (i) the continuous Maxwell--Boltzmann distribution and (ii) the
standard second-order Hermite-truncated discrete equilibrium. We show that
the two frameworks agree \emph{exactly} at second order, whereas at third
order the discrete equilibrium omits precisely $\rho\,u_\alpha u_\beta
u_\gamma$. The algebraic origin of the omission is identified as the parity of
the highest lattice moment entering the quadratic part of the equilibrium.
We then analyse the D3Q27 lattice component by component and prove that the
fully diagonal third-order moment $\Pi_{\alpha\alpha\alpha}$ is locked to
$\rho u_\alpha$ by the cubic identity $c_{i\alpha}^{3}=c_{i\alpha}$,
independently of the equilibrium; consequently the diagonal cubic term cannot
be restored by \emph{any} Hermite extension on a three-speed lattice, while
the mixed components are recoverable. The result is connected to the
$O(\mathrm{Ma}^{3})$ Galilean-invariance error of standard LB and to its
resolution through the cumulant collision operator, whose equilibrium
possesses only a finite number of non-vanishing cumulants.
\end{abstract}

\begin{keyword}
lattice Boltzmann method \sep equilibrium distribution \sep velocity moments
\sep Hermite expansion \sep Galilean invariance \sep cumulant collision
operator \sep D3Q27
\end{keyword}

\end{frontmatter}

\section{Introduction}
\label{sec:intro}

The lattice Boltzmann (LB) method recovers the Navier--Stokes equations
through the velocity moments of a discrete equilibrium distribution
$f_i^{\eq}$ defined on a finite velocity set $\{\mathbf{c}_i\}_{i=0}^{Q-1}$
with weights $\{w_i\}$ \citep{Qian1992,HeLuo1997,Kruger2017}. Because the
hydrodynamic stress is built from the second- and third-order moments of
$f_i^{\eq}$, the precise form of these moments determines which macroscopic
equations the scheme realises and at what order in the Mach number.

A recurring source of confusion is that the third-order moment is reported in
two different forms. Derivations that begin from continuous kinetic theory
(e.g.\ \citealp{HeLuo1997,ShanYuanChen2006}) list the third moment of the
\emph{full} Maxwell--Boltzmann distribution, which contains the cubic term
$\rho\,u_\alpha u_\beta u_\gamma$. Treatments that report the moment actually
realised by the standard second-order discrete equilibrium (e.g.\ the
textbook account of \citealp{Kruger2017}) omit this term. The two statements
are not in conflict: they describe different objects. The purpose of this
note is to make that distinction completely explicit by deriving all four
quantities --- the second- and third-order moments of the continuous
Maxwellian and of the discrete equilibrium --- and then to examine what the
D3Q27 lattice can and cannot represent.

The note is organised as follows. Section~\ref{sec:prelim} fixes the
continuous and discrete frameworks and states the two sets of moment
identities that every subsequent calculation reduces to.
Section~\ref{sec:second} treats the second-order moment and establishes exact
agreement between the two frameworks. Section~\ref{sec:third} treats the
third-order moment and isolates the cubic discrepancy.
Section~\ref{sec:d3q27} analyses the D3Q27 lattice component by component.
Section~\ref{sec:concl} collects the consequences for Galilean invariance and
for the cumulant collision operator.

\section{Preliminaries}
\label{sec:prelim}

\subsection{Continuous framework}

The continuous Maxwell--Boltzmann equilibrium in $D$ dimensions is
\begin{equation}
f^{\eq}(\boldsymbol{\xi})
=\frac{\rho}{(2\pi \cs^{2})^{D/2}}
\exp\!\left(-\frac{|\boldsymbol{\xi}-\mathbf{u}|^{2}}{2\cs^{2}}\right),
\label{eq:maxwell}
\end{equation}
where $\boldsymbol{\xi}$ is the microscopic velocity, $\rho$ the density,
$\mathbf{u}$ the macroscopic velocity and $\cs^{2}=RT$ the isothermal sound
speed squared. Define the ensemble average
$\langle\phi\rangle\equiv\int \phi\, f^{\eq}\dd^{D}\xi$ and the peculiar
velocity $\mathbf{C}=\boldsymbol{\xi}-\mathbf{u}$, so that under $\mathbf{C}$
the distribution \eqref{eq:maxwell} is a zero-mean isotropic Gaussian of
variance $\cs^{2}$ per component.

\begin{lemma}[Gaussian central moments]
\label{lem:gauss}
With the conventions above,
\begin{align}
\langle 1\rangle &= \rho, &
\langle C_\alpha\rangle &= 0, &
\langle C_\alpha C_\beta\rangle &= \rho \cs^{2}\delta_{\alpha\beta},
\label{eq:gauss012}\\
\langle C_\alpha C_\beta C_\gamma\rangle &= 0, &
\langle C_\alpha C_\beta C_\gamma C_\delta\rangle
&= \rho \cs^{4}\,\Deltafour, & &
\label{eq:gauss34}
\end{align}
with $\Deltafour=\delta_{\alpha\beta}\delta_{\gamma\delta}
+\delta_{\alpha\gamma}\delta_{\beta\delta}
+\delta_{\alpha\delta}\delta_{\beta\gamma}$; all odd-order central moments
vanish.
\end{lemma}

\begin{proof}
The Gaussian factorises over components, so a multi-index central moment
equals $\rho\prod_{\mu}I_{n_\mu}$, where $n_\mu$ is the number of repetitions
of direction $\mu$ and $I_{n}=\int x^{n}(2\pi\cs^{2})^{-1/2}
e^{-x^{2}/2\cs^{2}}\dd x$ with $I_0=1$, $I_1=0$, $I_2=\cs^{2}$, $I_3=0$,
$I_4=3\cs^{4}$. Any factor with odd $n_\mu$ annihilates the product, giving
the stated odd-order vanishing and \eqref{eq:gauss012}. For the fourth-order
moment the surviving index patterns are the three pairings, each contributing
$I_2 I_2=\cs^{4}$, which is the Isserlis/Wick result $\rho\cs^{4}\Deltafour$.
\end{proof}

\subsection{Discrete framework}

The standard LB equilibrium is the second-order Hermite truncation of
\eqref{eq:maxwell} evaluated on the lattice
\citep{Qian1992,HeLuo1997,ShanYuanChen2006},
\begin{equation}
f_i^{\eq}
= w_i\,\rho\left[
1
+\frac{c_{i\mu}u_\mu}{\cs^{2}}
+\frac{u_\mu u_\nu\big(c_{i\mu}c_{i\nu}-\cs^{2}\delta_{\mu\nu}\big)}{2\cs^{4}}
\right].
\label{eq:feq2}
\end{equation}
Standard isothermal velocity sets (D2Q9, D3Q15, D3Q19, D3Q27) are
Gauss--Hermite quadratures exact to polynomial degree five
\citep{HeLuo1997,ShanYuanChen2006}, so the weights satisfy the moment
identities
\begin{subequations}
\label{eq:lattice}
\begin{align}
\sumi w_i &= 1, &
\sumi w_i c_{i\alpha} &= 0,
\label{eq:M01}\\
\sumi w_i c_{i\alpha}c_{i\beta} &= \cs^{2}\delta_{\alpha\beta},  &
\sumi w_i c_{i\alpha}c_{i\beta}c_{i\gamma} &= 0,
\label{eq:M23}\\
\sumi w_i c_{i\alpha}c_{i\beta}c_{i\gamma}c_{i\delta}
&= \cs^{4}\,\Deltafour, &
\sumi w_i c_{i\alpha}c_{i\beta}c_{i\gamma}c_{i\delta}c_{i\epsilon} &= 0.
\label{eq:M45}
\end{align}
\end{subequations}
Equations \eqref{eq:M01}--\eqref{eq:M45} are labelled $M_0,\dots,M_5$ below.
The macroscopic moments are $\sum_i f_i^{\eq}=\rho$ and
$\sum_i f_i^{\eq}c_{i\alpha}=\rho u_\alpha$ by construction.

The third moment of the third-order equilibrium treated in
Section~\ref{sec:third-order-eq} requires one further identity, the sixth-order
weight moment, which is the \emph{first} member of the hierarchy that departs
from full isotropy. For the tensor-product lattices it reads
\begin{equation}
M^{(6)}_{\alpha\beta\gamma\delta\epsilon\zeta}
\equiv\sumi w_i\,c_{i\alpha}c_{i\beta}c_{i\gamma}c_{i\delta}c_{i\epsilon}c_{i\zeta}
=\cs^{6}\,\Delta^{(6)}_{\alpha\beta\gamma\delta\epsilon\zeta}
+A\,\delta^{[6]}_{\alpha\beta\gamma\delta\epsilon\zeta},
\label{eq:M6}
\end{equation}
where $\Delta^{(6)}$ is the fully symmetric sum of the $5!!=15$ pairings of six
indices into three Kronecker deltas, the diagonal indicator
$\delta^{[6]}_{\alpha\beta\gamma\delta\epsilon\zeta}$ equals unity iff all six
indices coincide and zero otherwise, and the anisotropic defect amplitude is
$A=m_6-15\cs^{6}$ with $m_6=\sumi w_i c_{ix}^{6}$ the one-dimensional sixth
moment. On D3Q27 the cubic identity $c_{ix}^{6}=c_{ix}^{2}$ gives
$m_6=\cs^{2}=\tfrac13$, against the isotropic value $15\cs^{6}=\tfrac59$, so
$A=-\tfrac29$. The fourth-order moment, by contrast, is fully isotropic
($m_4=3\cs^{4}$, consistent with \eqref{eq:M45}); thus $M_6$ is the lowest even
moment to break isotropy, and the breakage is confined to the fully diagonal
entries $M^{(6)}_{xxxxxx}=m_6$.

\section{The second-order moment}
\label{sec:second}

\begin{proposition}[Continuous second moment]
\label{prop:cont2}
$\displaystyle \int f^{\eq}\xi_\alpha\xi_\beta\dd^{D}\xi
=\rho u_\alpha u_\beta+\rho\cs^{2}\delta_{\alpha\beta}.$
\end{proposition}

\begin{proof}
Writing $\xi_\alpha=u_\alpha+C_\alpha$ and using Lemma~\ref{lem:gauss},
\begin{equation*}
\langle\xi_\alpha\xi_\beta\rangle
=u_\alpha u_\beta\langle 1\rangle
+u_\alpha\langle C_\beta\rangle
+u_\beta\langle C_\alpha\rangle
+\langle C_\alpha C_\beta\rangle
=\rho u_\alpha u_\beta+\rho \cs^{2}\delta_{\alpha\beta}. \qedhere
\end{equation*}
\end{proof}

\begin{proposition}[Discrete second moment]
\label{prop:disc2}
$\displaystyle \sum_i f_i^{\eq}c_{i\alpha}c_{i\beta}
=\rho\cs^{2}\delta_{\alpha\beta}+\rho u_\alpha u_\beta.$
\end{proposition}

\begin{proof}
Insert \eqref{eq:feq2} and split into the three terms of the bracket.
\emph{Constant term:} by $M_2$,
$\rho\sum_i w_i c_{i\alpha}c_{i\beta}=\rho\cs^{2}\delta_{\alpha\beta}$.
\emph{Linear term:} by $M_3$,
$(\rho u_\mu/\cs^{2})\sum_i w_i c_{i\mu}c_{i\alpha}c_{i\beta}=0$.
\emph{Quadratic term:} by $M_4$ and $M_2$,
\begin{align*}
\frac{\rho u_\mu u_\nu}{2\cs^{4}}
\Big[\sumi w_i c_{i\mu}c_{i\nu}c_{i\alpha}c_{i\beta}
-\cs^{2}\delta_{\mu\nu}\sumi w_i c_{i\alpha}c_{i\beta}\Big]
&=\frac{\rho u_\mu u_\nu}{2}
\big[\delta_{\mu\nu}\delta_{\alpha\beta}
+\delta_{\mu\alpha}\delta_{\nu\beta}
+\delta_{\mu\beta}\delta_{\nu\alpha}
-\delta_{\mu\nu}\delta_{\alpha\beta}\big]\\
&=\tfrac{\rho}{2}\big(u_\alpha u_\beta+u_\beta u_\alpha\big)
=\rho u_\alpha u_\beta.
\end{align*}
Summation gives the stated result.
\end{proof}

\begin{remark}
Propositions~\ref{prop:cont2} and \ref{prop:disc2} coincide: the second-order
moment, i.e.\ the momentum-flux tensor
$\Pi_{\alpha\beta}=\rho u_\alpha u_\beta+p\,\delta_{\alpha\beta}$ with
$p=\rho\cs^{2}$, is reproduced \emph{exactly} by the second-order discrete
equilibrium. This exactness is what allows standard LB to recover the viscous
stress correctly at leading order \citep{Kruger2017}.
\end{remark}

\section{The third-order moment}
\label{sec:third}

\begin{proposition}[Continuous third moment]
\label{prop:cont3}
\begin{equation}
\int f^{\eq}\xi_\alpha\xi_\beta\xi_\gamma\dd^{D}\xi
=\rho\cs^{2}\big(u_\alpha\delta_{\beta\gamma}
+u_\beta\delta_{\alpha\gamma}
+u_\gamma\delta_{\alpha\beta}\big)
+\rho\,u_\alpha u_\beta u_\gamma.
\label{eq:cont3}
\end{equation}
\end{proposition}

\begin{proof}
Expand $\xi_\alpha\xi_\beta\xi_\gamma
=(u_\alpha+C_\alpha)(u_\beta+C_\beta)(u_\gamma+C_\gamma)$ and average term by
term with Lemma~\ref{lem:gauss}. The triple-mean term yields
$\rho\,u_\alpha u_\beta u_\gamma$ (it survives because $\langle1\rangle=\rho$);
the three terms linear in $\mathbf{C}$ vanish through $\langle C\rangle=0$; the
three terms with two peculiar factors give
$u_\alpha\langle C_\beta C_\gamma\rangle+\text{(perm.)}
=\rho\cs^{2}(u_\alpha\delta_{\beta\gamma}+u_\beta\delta_{\alpha\gamma}
+u_\gamma\delta_{\alpha\beta})$; and the triple-peculiar term vanishes as an
odd Gaussian moment.
\end{proof}

\begin{proposition}[Discrete third moment]
\label{prop:disc3}
\begin{equation}
\sum_i f_i^{\eq}c_{i\alpha}c_{i\beta}c_{i\gamma}
=\rho\cs^{2}\big(u_\alpha\delta_{\beta\gamma}
+u_\beta\delta_{\alpha\gamma}
+u_\gamma\delta_{\alpha\beta}\big).
\label{eq:disc3}
\end{equation}
\end{proposition}

\begin{proof}
Insert \eqref{eq:feq2}. \emph{Constant term:} by $M_3$,
$\rho\sum_i w_i c_{i\alpha}c_{i\beta}c_{i\gamma}=0$.
\emph{Linear term:} by $M_4$,
\begin{equation*}
\frac{\rho u_\mu}{\cs^{2}}\sumi w_i
c_{i\mu}c_{i\alpha}c_{i\beta}c_{i\gamma}
=\frac{\rho u_\mu}{\cs^{2}}\cs^{4}
\big(\delta_{\mu\alpha}\delta_{\beta\gamma}
+\delta_{\mu\beta}\delta_{\alpha\gamma}
+\delta_{\mu\gamma}\delta_{\alpha\beta}\big)
=\rho\cs^{2}\big(u_\alpha\delta_{\beta\gamma}
+u_\beta\delta_{\alpha\gamma}+u_\gamma\delta_{\alpha\beta}\big).
\end{equation*}
\emph{Quadratic term:} the leading contraction is the fifth-order lattice
moment, which vanishes by $M_5$, and the subtracted piece vanishes by $M_3$;
hence the quadratic term contributes zero. Summation gives \eqref{eq:disc3}.
\end{proof}

\begin{remark}[The cubic discrepancy]
\label{rem:cubic}
Subtracting \eqref{eq:disc3} from \eqref{eq:cont3},
\begin{equation}
\underbrace{\int f^{\eq}\xi_\alpha\xi_\beta\xi_\gamma\dd^{D}\xi}_{\text{continuous}}
-\underbrace{\sum_i f_i^{\eq}c_{i\alpha}c_{i\beta}c_{i\gamma}}_{\text{discrete }O(u^{2})}
=\rho\,u_\alpha u_\beta u_\gamma .
\label{eq:gap}
\end{equation}
The structural reason for the omission is a parity argument visible in the
two proofs. In the second moment (Proposition~\ref{prop:disc2}) the highest
lattice moment reached by the quadratic part of the equilibrium is the
\emph{even} $M_4$, which is non-zero and supplies $\rho u_\alpha u_\beta$. In
the third moment the corresponding contraction reaches the \emph{odd} $M_5$,
which vanishes; the cubic term therefore has no source in a second-order
equilibrium. This missing contribution is the well-known nonlinear
($O(\mathrm{Ma}^{3})$) deviation of the lattice BGK models from the
Navier--Stokes stress \citep{QianOrszag1993,Kruger2017}.
\end{remark}

\section{The third-order equilibrium and its discrete third moment}
\label{sec:third-order-eq}

Remark~\ref{rem:cubic} located the cubic discrepancy in the truncation of the
equilibrium at second order. The natural remedy is to retain the next term of
the Hermite expansion of \eqref{eq:maxwell}. We write this third-order
equilibrium explicitly, then recompute the discrete third moment using the
sixth-order identity \eqref{eq:M6}. The cubic term is recovered for every mixed
index triple but not for the fully diagonal ones; the latter failure is shown,
in Section~\ref{sec:d3q27}, to be a property of the velocity set independent of
the equilibrium.

\subsection{The third-order Hermite equilibrium}
\label{sec:feq3}

The expansion of the Maxwellian \eqref{eq:maxwell} retained to third Hermite
order is \citep{Grad1949,HeLuo1997,ShanYuanChen2006}
\begin{equation}
f_i^{\eq,(3)}
= w_i\rho\left[
1
+\frac{c_{i\mu}u_\mu}{\cs^{2}}
+\frac{u_\mu u_\nu\,\mathcal{H}^{(2)}_{i,\mu\nu}}{2\cs^{4}}
+\frac{u_\mu u_\nu u_\lambda\,\mathcal{H}^{(3)}_{i,\mu\nu\lambda}}{6\cs^{6}}
\right],
\label{eq:feq3}
\end{equation}
with the second- and third-order tensor Hermite polynomials
\begin{align}
\mathcal{H}^{(2)}_{i,\mu\nu}&=c_{i\mu}c_{i\nu}-\cs^{2}\delta_{\mu\nu},
\label{eq:H2}\\
\mathcal{H}^{(3)}_{i,\mu\nu\lambda}&=c_{i\mu}c_{i\nu}c_{i\lambda}
-\cs^{2}\big(c_{i\mu}\delta_{\nu\lambda}+c_{i\nu}\delta_{\mu\lambda}
+c_{i\lambda}\delta_{\mu\nu}\big).
\label{eq:H3}
\end{align}
By index class, the ten independent components of \eqref{eq:H3} in three
dimensions are
\begin{equation}
\underbrace{\mathcal{H}^{(3)}_{xxx}=c_x^{3}-3\cs^{2}c_x}_{\text{diagonal}},
\qquad
\underbrace{\mathcal{H}^{(3)}_{xxy}=(c_x^{2}-\cs^{2})\,c_y}_{\text{two equal}},
\qquad
\underbrace{\mathcal{H}^{(3)}_{xyz}=c_xc_yc_z}_{\text{all distinct}},
\label{eq:H3components}
\end{equation}
together with their index permutations (three diagonal, six two-equal, one
all-distinct).

\begin{remark}[Vanishing of the diagonal Hermite mode on D3Q27]
\label{rem:H3vanish}
On D3Q27 the cubic identity $c_{ix}^{3}=c_{ix}$ together with $3\cs^{2}=1$ give
\begin{equation}
\mathcal{H}^{(3)}_{i,xxx}\big|_{\text{D3Q27}}
=c_{ix}^{3}-3\cs^{2}c_{ix}=c_{ix}-c_{ix}=0
\label{eq:H3xxxzero}
\end{equation}
identically over the lattice. The third-order Hermite subspace represented by
D3Q27 thus contains only the seven mixed modes; the three diagonal modes are the
zero vector. No equilibrium can populate a basis direction absent from the
lattice.
\end{remark}

\subsection{Discrete third moment of the third-order equilibrium}
\label{sec:disc3-order3}

\begin{proposition}[Discrete third moment under $f_i^{\eq,(3)}$]
\label{prop:disc3order3}
On D3Q27, the third moment of the third-order equilibrium \eqref{eq:feq3} is
\begin{equation}
\sum_i f_i^{\eq,(3)}c_{i\alpha}c_{i\beta}c_{i\gamma}
=\rho\cs^{2}\big(u_\alpha\delta_{\beta\gamma}+u_\beta\delta_{\alpha\gamma}
+u_\gamma\delta_{\alpha\beta}\big)
+\rho\,u_\alpha u_\beta u_\gamma\,\big(1-\delta^{[3]}_{\alpha\beta\gamma}\big),
\label{eq:disc3order3}
\end{equation}
where $\delta^{[3]}_{\alpha\beta\gamma}=1$ if $\alpha=\beta=\gamma$ and zero
otherwise. The cubic term is recovered exactly for every mixed index triple and
is absent for the fully diagonal components.
\end{proposition}

\begin{proof}
Write $\Pi_{\alpha\beta\gamma}=\Pi^{(2)}_{\alpha\beta\gamma}
+\Delta\Pi_{\alpha\beta\gamma}$, where $\Pi^{(2)}$ collects the first three terms
of \eqref{eq:feq3} and equals
$\rho\cs^{2}(u_\alpha\delta_{\beta\gamma}+u_\beta\delta_{\alpha\gamma}
+u_\gamma\delta_{\alpha\beta})$ by Proposition~\ref{prop:disc3}, and the
third-order term contributes
\begin{equation*}
\Delta\Pi_{\alpha\beta\gamma}
=\frac{\rho}{6\cs^{6}}\,u_\mu u_\nu u_\lambda\,
\Theta_{\mu\nu\lambda,\alpha\beta\gamma},
\qquad
\Theta_{\mu\nu\lambda,\alpha\beta\gamma}
=\sumi w_i\,\mathcal{H}^{(3)}_{i,\mu\nu\lambda}\,
c_{i\alpha}c_{i\beta}c_{i\gamma}.
\end{equation*}
The inverse Hermite relation $c_{i\alpha}c_{i\beta}c_{i\gamma}
=\mathcal{H}^{(3)}_{i,\alpha\beta\gamma}
+\cs^{2}(\delta_{\alpha\beta}c_{i\gamma}+\delta_{\alpha\gamma}c_{i\beta}
+\delta_{\beta\gamma}c_{i\alpha})$ splits $\Theta$ into a product of two
third-order Hermite polynomials plus cross terms
$\sumi w_i\,\mathcal{H}^{(3)}_{i,\mu\nu\lambda}c_{i\gamma}$. The cross terms are
degree-four polynomials, integrated exactly by the fifth-degree quadrature, and
therefore equal their continuum value
$\langle\mathcal{H}^{(3)}\mathcal{H}^{(1)}\rangle=0$ by Hermite orthogonality.
Hence
\begin{equation*}
\Theta_{\mu\nu\lambda,\alpha\beta\gamma}
=\sumi w_i\,\mathcal{H}^{(3)}_{i,\mu\nu\lambda}
\mathcal{H}^{(3)}_{i,\alpha\beta\gamma}.
\end{equation*}
This degree-six sum equals its continuum value
$\cs^{6}P_{\mu\nu\lambda,\alpha\beta\gamma}$, with
$P=\sum_{\sigma\in S_3}\delta_{\mu\sigma_1}\delta_{\nu\sigma_2}
\delta_{\lambda\sigma_3}$ the symmetric pairing of $(\mu\nu\lambda)$ with
$(\alpha\beta\gamma)$, provided no monomial $c_{ix}^{6}$ arises --- the only
entry of \eqref{eq:M6} that carries the defect $A$. Such a monomial requires
both factors to supply $c_{ix}^{3}$, that is $\mu=\nu=\lambda=x$ \emph{and}
$\alpha=\beta=\gamma=x$.

\emph{Mixed indices} ($\alpha,\beta,\gamma$ not all equal):
$\mathcal{H}^{(3)}_{\alpha\beta\gamma}$ carries no $c_{ix}^{3}$ term, so no
$c_{ix}^{6}$ arises for any $(\mu\nu\lambda)$ and the sum is exact. Contracting
with $u_\mu u_\nu u_\lambda$, each permutation in $P$ contributes
$u_\alpha u_\beta u_\gamma$, so
$u_\mu u_\nu u_\lambda\Theta=6\cs^{6}\,u_\alpha u_\beta u_\gamma$ and
$\Delta\Pi_{\alpha\beta\gamma}=\rho\,u_\alpha u_\beta u_\gamma$.

\emph{Diagonal indices} ($\alpha=\beta=\gamma=x$): by
Remark~\ref{rem:H3vanish}, $\mathcal{H}^{(3)}_{i,xxx}\equiv0$ over D3Q27, so the
correction associated with this mode in \eqref{eq:feq3} is identically zero and
contributes nothing, giving $\Delta\Pi_{xxx}=0$.
Section~\ref{sec:d3q27} reaches the same conclusion through the quadrature defect
$A$ of \eqref{eq:M6} and through an equilibrium-independent identity. Collecting
both cases yields \eqref{eq:disc3order3}.
\end{proof}

A concrete check of the mixed case is instructive. For $\Pi_{xyz}$ with
$x,y,z$ distinct, the second-order part vanishes; on the tensor-product lattice
the only non-vanishing internal indices are the permutations of $(x,y,z)$, each
yielding $\sumi w_i c_{ix}^{2}c_{iy}^{2}c_{iz}^{2}=m_2^{3}=\cs^{6}$. Summing the
six permutations,
$\Delta\Pi_{xyz}=(\rho/6\cs^{6})\cdot 6\cs^{6}\,u_xu_yu_z=\rho\,u_xu_yu_z$, in
agreement with the continuum value \eqref{eq:cont3}.

\section{Component-wise reality on the D3Q27 lattice}
\label{sec:d3q27}

Proposition~\ref{prop:disc3order3} established the recovery pattern by direct
computation with the third-order equilibrium. We now give an
equilibrium-independent argument that reaches the same conclusion for the
diagonal components, exposing it as a property of the velocity set alone, and
collect the result in a single table. Throughout this section we use the D3Q27
lattice in lattice units, $c_{i\alpha}\in\{-1,0,1\}$, $\cs^{2}=1/3$.

\subsection{The diagonal lock}

\begin{proposition}[Diagonal cubic lock]
\label{prop:lock}
On any lattice with $c_{i\alpha}\in\{-1,0,1\}$ (D2Q9, D3Q15, D3Q19, D3Q27),
\begin{equation}
\Pi_{\alpha\alpha\alpha}\equiv\sum_i f_i\,c_{i\alpha}^{3}=\rho u_\alpha
\qquad\text{for every }f_i,
\label{eq:lock}
\end{equation}
in particular independently of the choice or Hermite order of the
equilibrium.
\end{proposition}

\begin{proof}
Since $c_{i\alpha}\in\{-1,0,1\}$, the cubic identity $c_{i\alpha}^{3}
=c_{i\alpha}$ holds pointwise. Hence
$\sum_i f_i c_{i\alpha}^{3}=\sum_i f_i c_{i\alpha}=\rho u_\alpha$ for any
distribution.
\end{proof}

Because the continuous target \eqref{eq:cont3} requires
$\Pi_{\alpha\alpha\alpha}=3\rho\cs^{2}u_\alpha+\rho u_\alpha^{3}
=\rho u_\alpha+\rho u_\alpha^{3}$ (using $3\cs^{2}=1$), the diagonal cubic
term $\rho u_\alpha^{3}$ has no degree of freedom on a three-speed lattice and
cannot be represented.

\subsection{The sixth-order moment defect}

The same obstruction appears, equivalently, as a quadrature defect.
Augmenting \eqref{eq:feq2} with the third-order Hermite term and evaluating
its contribution to the diagonal moment $\Pi_{xxx}$ reduces, on the
one-dimensional factor D1Q3, to
\begin{equation}
\frac{u_x^{3}}{6\cs^{6}}\big(M_6-3\cs^{2}M_4\big),
\qquad
M_6=\sumi w_i c_{ix}^{6},\quad M_4=\sumi w_i c_{ix}^{4}.
\label{eq:sixth}
\end{equation}
For D1Q3 one has $M_4=\tfrac13$ (which equals $3\cs^{4}$, so the
fourth-order identity in \eqref{eq:M45} holds), whereas
$M_6=\tfrac13\neq 15\cs^{6}=\tfrac59$. Substituting,
$M_6-3\cs^{2}M_4=\tfrac13-3\cdot\tfrac13\cdot\tfrac13=0$, so the third-order
Hermite term contributes \emph{exactly zero} to the diagonal moment: the
lattice sixth-order defect cancels the very term it was introduced to
supply. This is the quadrature-level statement of Proposition~\ref{prop:lock},
and is consistent with the lattice's fifth-degree quadrature accuracy
\citep{ShanYuanChen2006}.

\subsection{Recoverable mixed components}

For a component in which no single direction carries a power exceeding four,
the relevant lattice moment is a tensor product of one-dimensional factors,
each exact, so the third-order Hermite term \emph{does} restore the cubic
contribution. For example $\sum_i w_i c_{ix}^{2}c_{iy}^{2}c_{iz}^{2}
=\cs^{6}$ matches the continuum, and the mixed moments are recovered:
\begin{align}
\Pi_{xyz}^{\text{cont}}&=\rho\,u_x u_y u_z, &
\Pi_{xxy}^{\text{cont}}&=\rho\cs^{2}u_y+\rho\,u_x^{2}u_y ,
\label{eq:mixed}
\end{align}
both attainable with the third-order Hermite equilibrium. Only the fully
diagonal components remain locked. Table~\ref{tab:summary} collects the three
representative cases.

\begin{table}[t]
\centering
\caption{Third-order equilibrium moment on D3Q27 ($\cs^{2}=1/3$, $3\cs^{2}=1$),
per unit $\rho$, for the three index classes. The second-order column is
Proposition~\ref{prop:disc3}; the Hermite-3 column is the best attainable on a
three-speed lattice; the Maxwell column is the continuum target
\eqref{eq:cont3}.}
\label{tab:summary}
\begin{tabular}{lccc}
\toprule
Component & Continuous Maxwell & Discrete $O(u^{2})$ & Discrete $+$ Hermite-3 \\
\midrule
$\Pi_{xxx}$ (diagonal) & $u_x+u_x^{3}$ & $u_x$ & $u_x$ \;(locked) \\
$\Pi_{xxy}$ (two equal) & $\tfrac13 u_y+u_x^{2}u_y$ & $\tfrac13 u_y$ &
$\tfrac13 u_y+u_x^{2}u_y$ \\
$\Pi_{xyz}$ (all distinct) & $u_x u_y u_z$ & $0$ & $u_x u_y u_z$ \\
\bottomrule
\end{tabular}
\end{table}

\section{Consequences and concluding remarks}
\label{sec:concl}

The three rows of Table~\ref{tab:summary} establish a strict ordering of
representable third moments,
\begin{equation}
\underbrace{\rho\cs^{2}(\cdots)}_{O(u^{2})\text{ truncation}}
\ \subsetneq\
\underbrace{\rho\cs^{2}(\cdots)+\rho\,u_\alpha u_\beta u_\gamma\big|_{\text{mixed only}}}_{\text{Hermite-3 on D3Q27}}
\ \subsetneq\
\underbrace{\rho\cs^{2}(\cdots)+\rho\,u_\alpha u_\beta u_\gamma\big|_{\text{all}}}_{\text{continuous Maxwell}} .
\label{eq:hierarchy}
\end{equation}
The gap between the second and third members is the diagonal cubic term, which
the three-speed lattice cannot carry. This is precisely the lattice-induced
$O(\mathrm{Ma}^{3})$ violation of Galilean invariance of the standard scheme
\citep{QianOrszag1993,ChikatamarlaKarlin2006}; its complete removal requires
multi-speed (higher-quadrature) velocity sets such as those constructed from
fifth-degree Gauss--Hermite rules \citep{ChikatamarlaKarlin2009}.

It is instructive to see why the cumulant collision operator
\citep{Geier2015} sidesteps the issue rather than attempting to restore the
diagonal raw moment, which Proposition~\ref{prop:lock} shows to be
impossible. The two-sided Laplace transform of \eqref{eq:maxwell} is
$M(\mathbf{k})/\rho=\exp\!\big(k_\mu u_\mu+\tfrac12\cs^{2}k_\mu k_\mu\big)$,
so the cumulant generating function $\ln(M/\rho)$ is a second-degree
polynomial: the Maxwellian possesses only a first cumulant
$\kappa_{1,\alpha}=u_\alpha$ and a second cumulant
$\kappa_{2,\alpha\beta}=\cs^{2}\delta_{\alpha\beta}$, with all higher
cumulants identically zero \citep{Geier2015}. Through the moment--cumulant
relation the cubic term in \eqref{eq:cont3} is the reducible contribution
$\kappa_{1,\alpha}\kappa_{1,\beta}\kappa_{1,\gamma}
=\rho\,u_\alpha u_\beta u_\gamma$, which is frame-dependent rather than a
genuine third-order cumulant. The cumulant operator relaxes the
Galilean-invariant cumulants directly, setting the third-order cumulant to
its (vanishing) equilibrium value; it therefore neither requires nor
falsely advertises recovery of the diagonal raw moment, which accounts for
its improved Galilean behaviour at high Reynolds number relative to a
Hermite-extended BGK equilibrium on the same D3Q27 lattice.

In summary, the apparently conflicting third-order results in the literature
are two correct statements about two different objects: the continuum
Maxwellian carries $\rho\,u_\alpha u_\beta u_\gamma$
(Proposition~\ref{prop:cont3}), the second-order discrete equilibrium does
not (Proposition~\ref{prop:disc3}), and on the D3Q27 lattice the diagonal
portion of that term is structurally unrepresentable
(Proposition~\ref{prop:lock}). The second-order moment, by contrast, is exact
in both frameworks (Propositions~\ref{prop:cont2}--\ref{prop:disc2}), which is
why standard LB recovers the Navier--Stokes stress at the intended order.

\begin{thebibliography}{99}

\bibitem[Bhatnagar et al.(1954)]{BGK1954}
P.~L. Bhatnagar, E.~P. Gross, M.~Krook,
A model for collision processes in gases. I. Small amplitude processes in
charged and neutral one-component systems,
Phys. Rev. 94 (1954) 511--525.
\href{https://doi.org/10.1103/PhysRev.94.511}{doi:10.1103/PhysRev.94.511}.

\bibitem[Chikatamarla and Karlin(2006)]{ChikatamarlaKarlin2006}
S.~S. Chikatamarla, I.~V. Karlin,
Entropy and Galilean invariance of lattice Boltzmann theories,
Phys. Rev. Lett. 97 (19) (2006) 190601.
\href{https://doi.org/10.1103/PhysRevLett.97.190601}{doi:10.1103/PhysRevLett.97.190601}.

\bibitem[Chikatamarla and Karlin(2009)]{ChikatamarlaKarlin2009}
S.~S. Chikatamarla, I.~V. Karlin,
Lattices for the lattice Boltzmann method,
Phys. Rev. E 79 (4) (2009) 046701.
\href{https://doi.org/10.1103/PhysRevE.79.046701}{doi:10.1103/PhysRevE.79.046701}.

\bibitem[Geier et al.(2015)]{Geier2015}
M.~Geier, M.~Sch\"onherr, A.~Pasquali, M.~Krafczyk,
The cumulant lattice Boltzmann equation in three dimensions: Theory and
validation,
Comput. Math. Appl. 70 (4) (2015) 507--547.
\href{https://doi.org/10.1016/j.camwa.2015.05.001}{doi:10.1016/j.camwa.2015.05.001}.

\bibitem[Grad(1949)]{Grad1949}
H.~Grad,
On the kinetic theory of rarefied gases,
Commun. Pure Appl. Math. 2 (4) (1949) 331--407.
\href{https://doi.org/10.1002/cpa.3160020403}{doi:10.1002/cpa.3160020403}.

\bibitem[He and Luo(1997)]{HeLuo1997}
X.~He, L.-S. Luo,
Theory of the lattice Boltzmann method: From the Boltzmann equation to the
lattice Boltzmann equation,
Phys. Rev. E 56 (6) (1997) 6811--6817.
\href{https://doi.org/10.1103/PhysRevE.56.6811}{doi:10.1103/PhysRevE.56.6811}.

\bibitem[Kr\"uger et al.(2017)]{Kruger2017}
T.~Kr\"uger, H.~Kusumaatmaja, A.~Kuzmin, O.~Shardt, G.~Silva, E.~M. Viggen,
The Lattice Boltzmann Method: Principles and Practice,
Graduate Texts in Physics, Springer, Cham, 2017.
\href{https://doi.org/10.1007/978-3-319-44649-3}{doi:10.1007/978-3-319-44649-3}.

\bibitem[Qian and Orszag(1993)]{QianOrszag1993}
Y.-H. Qian, S.~A. Orszag,
Lattice BGK models for the Navier--Stokes equation: Nonlinear deviation in
compressible regimes,
Europhys. Lett. 21 (3) (1993) 255--259.
\href{https://doi.org/10.1209/0295-5075/21/3/001}{doi:10.1209/0295-5075/21/3/001}.

\bibitem[Qian et al.(1992)]{Qian1992}
Y.-H. Qian, D.~d'Humi\`eres, P.~Lallemand,
Lattice BGK models for Navier--Stokes equation,
Europhys. Lett. 17 (6) (1992) 479--484.
\href{https://doi.org/10.1209/0295-5075/17/6/001}{doi:10.1209/0295-5075/17/6/001}.

\bibitem[Shan et al.(2006)]{ShanYuanChen2006}
X.~Shan, X.-F. Yuan, H.~Chen,
Kinetic theory representation of hydrodynamics: a way beyond the
Navier--Stokes equation,
J. Fluid Mech. 550 (2006) 413--441.
\href{https://doi.org/10.1017/S0022112005008153}{doi:10.1017/S0022112005008153}.

\end{thebibliography}

\end{document}
